% Section 7.3, from lectures/L07/README.md: the objectives, the questions,
% the further reading, and what the next chapter does with all of it.


\section{Review}
\label{c7:sec:review}

\needspace{6\baselineskip}
\subsection*{What you should be able to do now}
After this chapter, you should be able to:
\begin{itemize}
  \item Draw the small-signal schematic of any stage in this book, and say what each step discarded.
  \item Write $r_e$ from a collector current without looking it up.
  \item Derive gain, input resistance and output resistance for a common-emitter stage.
  \item State the emitter factor, and both of the things it decides.
  \item Say which node the emitter factor's boost appears at, and why it is invisible at another.
  \item Explain why a current-mirror load raises the gain by two independent mechanisms.
  \item Transfer any result to a MOSFET by substituting $r_s$, and name the one result that does not transfer.
  \item Compute a Miller capacitance, and say what a cascode does about it.
\end{itemize}

\needspace{6\baselineskip}
\subsection*{Questions to test yourself}
You should be able to answer:
\begin{enumerate}
  \item $r_e$ is 26 ohm at 1 mA. Is it a resistor? What happens to it if the stage is switched off?
  \item A stage has $R_C = 10$ kilohm and $R_E = 234$ ohm. What is its gain, and what would it be with the emitter resistor bypassed?
  \item The emitter factor is ten. Name the two quantities it relates, and say which of them is a property of the transistor rather than of the stage.
  \item Why does adding an emitter resistor barely change the output resistance of a resistively loaded stage, when it multiplies the resistance looking into the collector by ten?
  \item A current mirror as a load raises the gain for two separate reasons. Name both.
  \item A stage with a gain of 385 and 4 pF of collector-base capacitance is driven from 1 kilohm. Where is its input pole?
  \item A cascode is often described as a new circuit. In what sense is it a stage you have already analysed?
\end{enumerate}

\needspace{6\baselineskip}
\subsection*{Further reading}
\begin{itemize}
  \item \emph{The Art of Electronics}, Horowitz and Hill, chapter 2, for the small-signal model, and \emph{Microelectronic Circuits}, Sedra and Smith, if you want the hybrid-pi form instead.
\end{itemize}

\needspace{6\baselineskip}
\subsection*{Looking ahead}
\begin{itemize}
  \item Chapter~\ref{c8:ch:follower} is the stage with no voltage gain at all, and why almost every design contains one. Followers, impedance transformation, and the output stage they turn into.
\end{itemize}
